\documentclass[11pt,aspectratio=169]{ctexbeamer}

\usetheme{default}
\setbeamertemplate{navigation symbols}{}

\usepackage{amsmath}
\usepackage{booktabs}

\definecolor{DeepBlue}{RGB}{25,74,122}
\definecolor{SoftInk}{RGB}{53,59,66}
\definecolor{PaperBg}{RGB}{249,248,245}
\definecolor{PanelBlue}{RGB}{236,243,249}

\setbeamercolor{normal text}{fg=SoftInk,bg=PaperBg}
\setbeamercolor{structure}{fg=DeepBlue}
\setbeamercolor{frametitle}{fg=DeepBlue,bg=PaperBg}
\setbeamercolor{title}{fg=DeepBlue}
\setbeamercolor{block title}{fg=DeepBlue,bg=PanelBlue}
\setbeamercolor{block body}{fg=SoftInk,bg=white}

\newcommand{\Core}[1]{
  \begin{block}{核心结论}
    #1
  \end{block}
}

\newcommand{\LectureMeta}[2]{
  \title{大模型推理系统与实践}
  \subtitle{#1}
  \author{张书豪}
  \institute{华中科技大学计算机学院}
  \date{2027 · #2}
}

\newcommand{\PaperRef}[3]{
  \item \textbf{#1}（#2）\hfill\href{#3}{原文链接}
}

\newcommand{\ResearchSources}[1]{
  \begin{frame}{本讲参考资料}
    \small
    \begin{itemize}
      #1
    \end{itemize}
    \vspace{0.3em}
    \textcolor{SoftInk}{课件中的论文数字均指作者在特定模型、硬件和工作负载下报告的结果，不代表无条件性能。}
  \end{frame}
}

\newcommand{\InternalCaseNote}{
  \small\textcolor{SoftInk}{说明：本页为课程研究案例或教学假设，不代表已发表论文结论；只有在实验包给出模型、硬件、工作负载和原始记录时，数字才可作为结果引用。}
}

% 14 讲细分版的统一参考资料入口。论文数字的解释边界见课程内容基线。

\newcommand{\LectureOneSources}{\ResearchSources{
  \PaperRef{vLLM: Easy, Fast, and Cheap LLM Serving with PagedAttention}{SOSP 2023}{https://doi.org/10.1145/3600006.3613165}
  \PaperRef{SGLang: Efficient Execution of Structured Language Model Programs}{arXiv:2312.07104}{https://arxiv.org/abs/2312.07104}
  \PaperRef{nano-vllm-hust 教学代码}{课程仓库}{https://github.com/intellistream/intro-to-llm-inference-engines}
}}

\newcommand{\LectureTwoSources}{\ResearchSources{
  \PaperRef{vLLM / PagedAttention}{SOSP 2023}{https://doi.org/10.1145/3600006.3613165}
  \PaperRef{Sarathi-Serve}{OSDI 2024}{https://www.usenix.org/conference/osdi24/presentation/agrawal}
  \PaperRef{DistServe}{OSDI 2024}{https://www.usenix.org/conference/osdi24/presentation/zhong-yinmin}
}}

\newcommand{\LectureThreeSources}{\ResearchSources{
  \PaperRef{vLLM / PagedAttention}{SOSP 2023}{https://doi.org/10.1145/3600006.3613165}
  \PaperRef{DistServe}{OSDI 2024}{https://www.usenix.org/conference/osdi24/presentation/zhong-yinmin}
  \PaperRef{Llumnix}{OSDI 2024}{https://www.usenix.org/conference/osdi24/presentation/sun-biao}
}}

\newcommand{\LectureFourSources}{\ResearchSources{
  \PaperRef{Orca}{OSDI 2022}{https://www.usenix.org/conference/osdi22/presentation/yu}
  \PaperRef{Sarathi-Serve}{OSDI 2024}{https://www.usenix.org/conference/osdi24/presentation/agrawal}
  \PaperRef{FastServe}{NSDI 2026}{https://www.usenix.org/conference/nsdi26/presentation/wu-bingyang}
}}

\newcommand{\LectureFiveSources}{\ResearchSources{
  \PaperRef{vLLM / PagedAttention}{SOSP 2023}{https://doi.org/10.1145/3600006.3613165}
  \PaperRef{vAttention}{arXiv:2405.04437}{https://arxiv.org/abs/2405.04437}
  \PaperRef{Rethinking KV Cache Compression}{MLSys 2025}{https://proceedings.mlsys.org/paper_files/paper/2025/hash/26289c647c6828e862e271ca3c490486-Abstract-Conference.html}
  \PaperRef{Mooncake}{FAST 2025}{https://www.usenix.org/conference/fast25/presentation/qin}
}}

\newcommand{\LectureSixSources}{\ResearchSources{
  \PaperRef{DistServe}{OSDI 2024}{https://www.usenix.org/conference/osdi24/presentation/zhong-yinmin}
  \PaperRef{RAGCache}{arXiv:2404.12457}{https://arxiv.org/abs/2404.12457}
  \PaperRef{Mooncake}{FAST 2025}{https://www.usenix.org/conference/fast25/presentation/qin}
  \PaperRef{NEO}{MLSys 2025}{https://proceedings.mlsys.org/paper_files/paper/2025/hash/66a026c0d17040889b50f0dfa650e5e0-Abstract-Conference.html}
}}

\newcommand{\LectureSevenSources}{\ResearchSources{
  \PaperRef{vLLM / PagedAttention}{SOSP 2023}{https://doi.org/10.1145/3600006.3613165}
  \PaperRef{SGLang}{arXiv:2312.07104}{https://arxiv.org/abs/2312.07104}
  \PaperRef{DistServe}{OSDI 2024}{https://www.usenix.org/conference/osdi24/presentation/zhong-yinmin}
  \PaperRef{Mooncake}{FAST 2025}{https://www.usenix.org/conference/fast25/presentation/qin}
}}

\newcommand{\LectureEightSources}{\ResearchSources{
  \PaperRef{FlashAttention}{NeurIPS 2022}{https://arxiv.org/abs/2205.14135}
  \PaperRef{FlexGen}{ICML 2023}{https://proceedings.mlr.press/v202/sheng23a.html}
  \PaperRef{PowerInfer}{SOSP 2024 / arXiv:2312.12456}{https://arxiv.org/abs/2312.12456}
  \PaperRef{FlashInfer}{MLSys 2025}{https://proceedings.mlsys.org/paper_files/paper/2025/hash/dbf02b21d77409a2db30e56866a8ab3a-Abstract-Conference.html}
}}

\newcommand{\LectureNineSources}{\ResearchSources{
  \PaperRef{DEEPSERVE}{USENIX ATC 2025}{https://www.usenix.org/conference/atc25/presentation/hu-junhao}
  \PaperRef{QServe}{MLSys 2025}{https://proceedings.mlsys.org/paper_files/paper/2025/hash/fbe2b2f74a2ece8070d8fb073717bda6-Abstract-Conference.html}
  \PaperRef{vLLM-HUST}{开源维护仓库}{https://github.com/vLLM-HUST/vllm-hust}
}}

\newcommand{\LectureTenSources}{\ResearchSources{
  \PaperRef{vLLM / PagedAttention}{SOSP 2023}{https://doi.org/10.1145/3600006.3613165}
  \PaperRef{Revisiting Disaggregated LLM Serving}{2025 preprint, arXiv:2601.08833}{https://arxiv.org/abs/2601.08833}
  \PaperRef{Rethinking KV Cache Compression}{MLSys 2025}{https://proceedings.mlsys.org/paper_files/paper/2025/hash/26289c647c6828e862e271ca3c490486-Abstract-Conference.html}
}}

\newcommand{\LectureElevenSources}{\ResearchSources{
  \PaperRef{MLPerf Inference}{MLCommons benchmark}{https://mlcommons.org/benchmarks/inference-datacenter/}
  \PaperRef{ACM Artifact Review and Badging}{ACM policy}{https://www.acm.org/publications/policies/artifact-review-and-badging-current}
  \PaperRef{Revisiting Disaggregated LLM Serving}{2025 preprint, arXiv:2601.08833}{https://arxiv.org/abs/2601.08833}
}}

\newcommand{\LectureTwelveSources}{\ResearchSources{
  \PaperRef{nano-vllm-hust}{课程教学代码}{https://github.com/intellistream/intro-to-llm-inference-engines}
  \PaperRef{vLLM-HUST}{HUST 维护的 vLLM fork}{https://github.com/vLLM-HUST/vllm-hust}
  \PaperRef{上游 vLLM}{开源项目}{https://github.com/vllm-project/vllm}
}}

\newcommand{\LectureThirteenSources}{\ResearchSources{
  \PaperRef{课程内容基线与实验材料}{课程仓库}{https://github.com/intellistream/intro-to-llm-inference-engines}
  \PaperRef{DEEPSERVE}{USENIX ATC 2025}{https://www.usenix.org/conference/atc25/presentation/hu-junhao}
  \PaperRef{GhostServe}{MLSys 2026}{https://proceedings.mlsys.org/paper_files/paper/2026/hash/d99e8e80a6c41e148db686918dd7eab3-Abstract-Conference.html}
}}

\newcommand{\LectureFourteenSources}{\ResearchSources{
  \PaperRef{课程完整材料与后续更新}{课程仓库}{https://github.com/intellistream/intro-to-llm-inference-engines}
  \PaperRef{SpecGen}{MLSys 2026}{https://proceedings.mlsys.org/paper_files/paper/2026/hash/66a026c0d17040889b50f0dfa650e5e0-Abstract-Conference.html}
  \PaperRef{GhostServe}{MLSys 2026}{https://proceedings.mlsys.org/paper_files/paper/2026/hash/d99e8e80a6c41e148db686918dd7eab3-Abstract-Conference.html}
}}



\LectureMeta{第 10 讲 论文比较方法与实验可信性}{研究生课程课件}

\begin{document}

\begin{frame}[plain]
  \titlepage
\end{frame}

\begin{frame}{教学目标}
  \Core{本讲训练系统论文阅读能力，重点在于从论文中提炼真实研究问题与比较维度。}
  \begin{itemize}
    \item 识别瓶颈假设和控制点变化。
    \item 分析机制代价与适用边界。
    \item 判断实验是否足以支撑结论。
  \end{itemize}
\end{frame}

\begin{frame}{系统论文的阅读重点}
  \Core{阅读系统论文时，首要任务是重建问题定义而非简单记忆图表。}
  \begin{itemize}
    \item 论文究竟重新定义了什么问题？
    \item 改变了哪个控制点或状态对象？
    \item 引入了哪些额外代价？
  \end{itemize}
\end{frame}

\begin{frame}{五段式精读框架}
  \Core{一个清晰的精读报告至少应包含五个部分。}
  \begin{enumerate}
    \item 问题假设
    \item 核心机制
    \item 额外代价
    \item 适用边界
    \item 实验支撑度
  \end{enumerate}
\end{frame}

\begin{frame}{什么叫“问题假设”}
  \Core{很多学生读论文只看作者做了什么，却忽略了作者默认哪些前提成立。}
  \begin{itemize}
    \item workload 是否是在线动态到达。
    \item 瓶颈是否真的在作者声称的部分。
    \item 硬件、模型和系统配置是否会改变结论。
  \end{itemize}
\end{frame}

\begin{frame}{忽略问题假设的风险}
  \Core{曲线和表格通常建立在特定假设、工作负载与平台条件之上，脱离前提容易导致误判。}
  \begin{itemize}
    \item 看似很大的加速可能只在窄场景成立。
    \item 看似中性的结果可能在更公平基线下并不成立。
    \item 论文真正的贡献常常在问题重构而不是数值本身。
  \end{itemize}
\end{frame}

\begin{frame}{如何识别真正的机制贡献}
  \Core{所谓“机制贡献”，应当表现为控制点、状态组织或执行路径被实质性改变。}
  \begin{itemize}
    \item 只是换了术语，不算机制。
    \item 只是重新组合已有做法，也未必构成实质贡献。
    \item 真正有价值的机制应当改变系统行为边界。
  \end{itemize}
\end{frame}

\begin{frame}{如何判断实验支撑度}
  \Core{实验的任务在于论证问题假设与机制收益是否成立。}
  \begin{itemize}
    \item workload 是否贴合目标场景。
    \item 对照组是否公平。
    \item 指标是否覆盖端到端效果与关键探针。
  \end{itemize}
\end{frame}

\begin{frame}{讲解：论文比较方法课真正训练的能力}
  \Core{这讲不是教学生复述论文摘要，而是训练把论文拆成“问题、机制、代价、边界、证据”的能力。}
  \begin{itemize}
    \item 先问作者到底重构了哪个系统问题，而不是先看图表大不大。
    \item 再问机制改变了什么控制点，为什么这个点值得改。
    \item 最后问实验到底支撑了哪些结论，还有哪些只是作者的延伸解释。
  \end{itemize}
\end{frame}

\begin{frame}{讲解：为什么系统论文不能只比“数字是否更高”}
  \Core{论文比较的核心不是找一个冠军，而是判断一篇工作是否真的把问题讲清、做实、证实。}
  \begin{itemize}
    \item 工程量很大，不等于研究判断就清楚。
    \item 加速很明显，不等于对照就公平、边界就被说明。
    \item 所以课程希望学生形成的不是“排名思维”，而是“证据与边界思维”。
  \end{itemize}
\end{frame}

\begin{frame}{例子：同一篇系统论文应该先看哪里（1）}
  \Core{论文比较方法课不是教学生背结论，而是教他们迅速识别论文的真正控制点和证据边界。}
\begin{block}{第一步：找问题定义}
    这篇论文到底在解决哪一个系统对象的什么痛点？
  \end{block}
\end{frame}

\begin{frame}{例子：同一篇系统论文应该先看哪里（2）}
  \Core{论文比较方法课不是教学生背结论，而是教他们迅速识别论文的真正控制点和证据边界。}
\begin{block}{第二步：找机制变化}
    它改了控制点、状态组织、执行路径，还是实验方法本身？
  \end{block}
\end{frame}

\begin{frame}{例子：同一篇系统论文应该先看哪里（3）}
  \Core{论文比较方法课不是教学生背结论，而是教他们迅速识别论文的真正控制点和证据边界。}
\begin{block}{第三步：找证据边界}
    它的实验到底支撑了哪部分结论，还有哪部分只是推测？
  \end{block}
\end{frame}

\begin{frame}{例子：把 vLLM 论文拆成五段会看到什么（1）}
  \Core{经典论文最好拿来做结构拆解，因为它能帮助学生建立统一的阅读框架。}
\begin{block}{假设}
    动态 serving 的关键瓶颈之一是 KV 管理方式不支持高效 continuous batching。
  \end{block}
\end{frame}

\begin{frame}{例子：把 vLLM 论文拆成五段会看到什么（2）}
  \Core{经典论文最好拿来做结构拆解，因为它能帮助学生建立统一的阅读框架。}
\begin{block}{机制}
    用 PagedAttention 改写状态组织方式，而不是只调度得更频繁。
  \end{block}
\end{frame}

\begin{frame}{例子：把 vLLM 论文拆成五段会看到什么（3）}
  \Core{经典论文最好拿来做结构拆解，因为它能帮助学生建立统一的阅读框架。}
\begin{block}{边界}
    论文证明了什么条件下收益成立，但不等于所有 serving 问题都因此解决。
  \end{block}
\end{frame}

\begin{frame}{代码例子：把论文机制映射回真实代码（1）}
  \Core{论文阅读真正有价值的时刻，是你能指出机制落到代码里大概对应哪类对象。}
\begin{block}{代码映射示例}
    {\ttfamily\small
    \detokenize{BlockManager / block_table -> KV organization}\par
    \detokenize{Scheduler -> control point}\par
    \detokenize{bench.py / scripts -> evidence path}\par
    }
  \end{block}
\end{frame}

\begin{frame}{代码例子：把论文机制映射回真实代码（2）}
  \Core{论文阅读真正有价值的时刻，是你能指出机制落到代码里大概对应哪类对象。}
\begin{block}{课堂应追问}
    如果一篇论文的核心机制在代码里找不到等价控制点，你该如何判断它是否被真实实现？
  \end{block}
\end{frame}

\begin{frame}{代码例子：把论文机制映射回真实代码（3）}
  \Core{论文阅读真正有价值的时刻，是你能指出机制落到代码里大概对应哪类对象。}
\begin{block}{课程结论}
    论文比较不能停在图表和术语层，必须回到对象、控制点和可验证实现。
  \end{block}
\end{frame}

\begin{frame}{代码例子：系统论文精读报告模板（1）}
  \Core{这类课程最好给学生一份“可直接套用但不偷懒”的读论文结构模板。}
\begin{block}{模板片段}
    {\ttfamily\small
    \detokenize{1. Problem and system object}\par
    \detokenize{2. Mechanism and changed control point}\par
    \detokenize{3. Cost / boundary / hidden assumptions}\par
    \detokenize{4. Evidence and what is still unsupported}\par
    }
  \end{block}
\end{frame}

\begin{frame}{代码例子：系统论文精读报告模板（2）}
  \Core{这类课程最好给学生一份“可直接套用但不偷懒”的读论文结构模板。}
\begin{block}{课堂应追问}
    为什么“论文摘要改写一遍”不能算合格精读报告？
  \end{block}
\end{frame}

\begin{frame}{代码例子：系统论文精读报告模板（3）}
  \Core{这类课程最好给学生一份“可直接套用但不偷懒”的读论文结构模板。}
\begin{block}{教学意义}
    模板的作用不是降低难度，而是强制学生用系统语言组织判断。
  \end{block}
\end{frame}

\begin{frame}{例子：什么叫 synthetic gain（1）}
  \Core{课程里必须训练学生识别“图很好看，但比较并不公平”的论文叙事。}
\begin{block}{迹象一}
    baseline 选择过弱，或者没有把现有系统最佳实践纳入比较。
  \end{block}
\end{frame}

\begin{frame}{例子：什么叫 synthetic gain（2）}
  \Core{课程里必须训练学生识别“图很好看，但比较并不公平”的论文叙事。}
\begin{block}{迹象二}
    workload 过于单一，恰好只适合提出的方法。
  \end{block}
\end{frame}

\begin{frame}{例子：什么叫 synthetic gain（3）}
  \Core{课程里必须训练学生识别“图很好看，但比较并不公平”的论文叙事。}
\begin{block}{迹象三}
    只给出端到端大图，不给局部探针与失败条件，导致因果链断裂。
  \end{block}
\end{frame}

\begin{frame}{研究案例：同一个结果可以支撑不同强度的 claim}
  \Core{系统论文写作不是把结果写满，而是把 claim 压到证据能够承受的范围。}
  \begin{itemize}
    \item KV-Delta 可说：repeated movement 下减少重复 KV image，并在 m 增大时改善 E2E。
    \item 不能直接说：所有迁移指标都改善，因为 TTFT / disruption 仍有反例。
    \item 课堂问题：如果只保留一句贡献，该写成 protocol win、transport win 还是 install-path win？
  \end{itemize}
\end{frame}

\begin{frame}{研究案例：负结果也可以成为系统论文材料}
  \Core{负结果不是失败汇报，而是帮助读者理解系统边界的证据。}
  \begin{itemize}
    \item Segment Reuse：能 commit 但输给 native prefix，说明 committed path 成本尚未压过强 baseline。
    \item Tiered KV：movement/restore 触发条件不足或代价过高，说明 admission policy 是核心问题。
    \item MLP Materialization：单算子占比过小，说明优化目标应转向边界与写回路径。
  \end{itemize}
\end{frame}

\begin{frame}{研究案例：Semantic MapReduce 的 claim 边界}
  \Core{SAGE 案例适合训练学生区分 orchestration claim、coverage claim 和 LLM quality claim。}
  \begin{itemize}
    \item 可以说：workflow / stream / semantic reduce 能组织大规模分析证据链。
    \item 需要谨慎：coverage 不足时 F1 下降，不能简单归因于 LLM reducer 不好。
    \item 不应声称：SAGE 已经替代 Spark/Flink/Ray 或已经证明所有 LLM reducer 稳定获胜。
  \end{itemize}
\end{frame}

\begin{frame}{例子：工程量大不等于研究贡献大（1）}
  \Core{系统论文比较方法课另一个关键目标，是帮助学生区分“做了很多”和“回答了好问题”。}
\begin{block}{工程量大}
    代码多、模块多、图表多，但问题定义可能仍然模糊。
  \end{block}
\end{frame}

\begin{frame}{例子：工程量大不等于研究贡献大（2）}
  \Core{系统论文比较方法课另一个关键目标，是帮助学生区分“做了很多”和“回答了好问题”。}
\begin{block}{研究贡献大}
    可能只改变一个关键控制点，却清楚说明了为什么这个点改变了系统边界。
  \end{block}
\end{frame}

\begin{frame}{例子：工程量大不等于研究贡献大（3）}
  \Core{系统论文比较方法课另一个关键目标，是帮助学生区分“做了很多”和“回答了好问题”。}
\begin{block}{课堂结论}
    论文比较首先比较问题与证据，其次才比较工程复杂度。
  \end{block}
\end{frame}

\begin{frame}{课堂练习：用五句话否定一篇论文（1）}
  \Core{真正的精读能力之一，是能用短而硬的判断指出论文的边界。}
\begin{block}{练习一}
    用一句话写出论文假设最脆弱的地方。
  \end{block}
\end{frame}

\begin{frame}{课堂练习：用五句话否定一篇论文（2）}
  \Core{真正的精读能力之一，是能用短而硬的判断指出论文的边界。}
\begin{block}{练习二}
    用一句话写出最需要补充的实验。
  \end{block}
\end{frame}

\begin{frame}{课堂练习：用五句话否定一篇论文（3）}
  \Core{真正的精读能力之一，是能用短而硬的判断指出论文的边界。}
\begin{block}{练习三}
    用一句话写出：如果这个实验结果反过来，整篇论文最先倒塌的结论是什么。
  \end{block}
\end{frame}


\begin{frame}{课堂练习}
  \Core{课堂练习可围绕代表论文的独立拆解展开，以训练结构化比较能力。}
  \begin{itemize}
    \item 每组选择一篇代表论文。
    \item 用五段式框架做 10 分钟汇报。
    \item 汇报中必须包含一个会否定作者结论的反例设想。
  \end{itemize}
\end{frame}

\begin{frame}{课堂讨论}
  \Core{系统论文阅读的重点在于判断问题是否真实、机制是否必要以及结果是否可信。}
  \begin{itemize}
    \item 如何区分“新机制”和“新命名”？
    \item 如何判断一个对照实验是否公平？
    \item 什么样的结果最应该要求作者额外解释？
  \end{itemize}
\end{frame}

\begin{frame}{本讲总结}
  \Core{规范的论文阅读能力能够直接提升研究选题质量与实验设计质量。}
  \begin{itemize}
    \item 论文是问题建模与机制设计的分析样本。
    \item 论文阅读的最终目标是形成稳定的研究判断能力。
  \end{itemize}
\end{frame}

\begin{frame}{对应 Tutorial}
  \Core{本讲对应 Tutorial 9。先独立作答，再对照下一页参考答案。}
  \begin{itemize}
    \item 文件：\texttt{\detokenize{build/tutorials/Tutorial_09_论文比较方法.pdf}}
    \item 动手运行、日志记录与提交要求统一查看 \texttt{\detokenize{build/experiments/}} 和对应 experiment 源目录下的 \texttt{\detokenize{assignment_spring-2026/}}。
  \end{itemize}
\end{frame}

\end{document}